% Copyright 2025 the authors. All rights reserved.

\documentclass{article}
\usepackage[utf8]{inputenc}
\usepackage[letterpaper]{geometry}
\usepackage{setspace}
\usepackage{amsmath, amsfonts, mathrsfs, bm}

% typesetting issues
\setstretch{1.08}
\addtolength{\topmargin}{-0.30in}
\addtolength{\textheight}{1.60in}
\setlength{\textwidth}{5.0in}
\setlength{\oddsidemargin}{0.5\paperwidth}\addtolength{\oddsidemargin}{-1.0in}\addtolength{\oddsidemargin}{-0.5\textwidth}
\pagestyle{myheadings}
\markboth{foo}{\sffamily Hogg \& Saydjari / Fitting a hyperplane with uncertainties in all directions}
\renewcommand{\newblock}{} % this adjusts the bibliography style.
\frenchspacing\sloppy\sloppypar\raggedbottom

% text macros
\renewcommand{\paragraph}[1]{\bigskip\par\noindent\textbf{#1}~---}
\newcommand{\documentname}{\textsl{Note}}
\newcommand{\sectionname}{Section}
\newcommand{\secref}[1]{\sectionname~\ref{#1}}
\renewcommand{\figurename}{Figure}
\newcommand{\figref}[1]{\figurename~\ref{#1}}
\newcommand{\foreign}[1]{\textsl{#1}}

% math macros
\renewcommand{\vec}[1]{\bm{#1}}
\newcommand{\mat}[1]{\mathsf{#1}}
\newcommand{\given}{\,|\,}

\title{\bfseries Fitting a hyperplane to data when there are uncertainties in all directions}
\author{Hogg \& Saydjari}
\date{May 2025}
\begin{document}
\maketitle\thispagestyle{empty}

The setup is as follows:
We have $N$ data points $\vec{x}_i\in\mathbb{R}^D$ ($1\leq i\leq N$).
Each of these has a $D\times D$ covariance matrix $\mat{C}_i$ representing (accurately) the $D$-dimensional variance tensor of the zero-mean additive Gaussian noise that has moved $\vec{x}_i$ away from its ``true'' value $\vec{\tilde{x}}_i$.
The assumption is that these ``true'' values lie on a $(D-1)$-dimensional hyperplane in the $D$ space.
In the case of $D=2$, this reduces to the problem of fitting a line to two-dimensional data, which have uncertainties in both dimensions (\cite{HBL}, Section~WHATEVER).

There are two approaches, both slightly annoying.
The first, and most consistent with the content of \cite{HBL}, is to go Bayesian.
In this case, you can think of the hyperplane as being a prior on the ``true'' locations $\vec{\tilde{x}}_i$ of each data point, and then the inference of the location and orientation of the hyperplane takes the form of a hierarchical inference.
The second, and requiring far fewer assumptions, is to go frequentist.
In this case, we find the maximum-log-likelihood location of each point on any proposed hyperplane, and add up those profile log likelihood values to get the log profile likelihood of the line.

In both of these cases, we have choices about parameterization of the hyperplane.
The option we choose here is that there plane is defined by a set of orthogonal unit vectors $\vec{\hat{u}}_k$ ($1\leq k\leq D$), such that any point on the plane can be described with a $(D-1)$-vector $\vec{s}$ of coordinates $s_k$ ($1\leq k\leq (D-1)$)
\begin{align}
    \vec{\tilde{x}}(\vec{s}) &= a\,\vec{\hat{u}}_D + \sum_{k=1}^{D-1} s_k\,\vec{\hat{u}}_k ~,
\end{align}
where $a$ gives the displacement of the hyperplane away from the origin in the perpendicular direction.
The goal of inference is to obtain a good estimate of the orientation of the hyperplane (the hyperspace spanned by the $\vec{\hat{u}}_k$ for $k<D$) and its position $a$.
This parameterization has degeneracies, which we could break if necessary;
fundamentally the orientation of the hyperplane is set by the orientation of its normal vector $\vec{\hat{u}}_D$.
This is all illustrated for $D=2$ in \figref{fig:setup}.

\paragraph{Bayesian approach}
The Bayesians need probability distributions over latents; we can provide.
We can imagine that the probability density for the $(D-1)$-dimensional position $\vec{s}$ on the hyperplane is a limit of a probability of being at a $D$-dimensional position $\vec{x}$ in the space.
For example, the probability density for points in the plane could be drawn from a Gaussian with mean $\mu$ and variance tensor $\mat{\Lambda}$ defined by
\begin{align}
    \mu &= a\,\vec{\hat{u}}_D \\
    \mat{\Lambda} &= \lim_{\epsilon\rightarrow 0}\left[ \epsilon^{(D-1)}\,\vec{\hat{u}}_D\,\vec{\hat{u}}_D^\top + \sum_{k=1}^{D-1} \frac{1}{\epsilon}\,\vec{\hat{u}}_k\,\vec{\hat{u}}_k^\top\right] ~,
\end{align}
where we have taken the limit of an infinitely thin, infinitely extended Gaussian, and implicitly the $D$-vectors $\vec{\hat{u}}_k$ are column vectors.
The power $(D-1)$ is employed to give this variance tensor a determinant of unity.

The Gaussian with mean $\mu$ and variance $\mat{\Lambda}$ is the probability distribution for the the ``true'' positions of any points that live on the hyperplane.
The noisy observed points $\vec{x}_i$ are displaced off of the plane by (zero-mean) Gaussian noise with known variance tensors $\mat{C}_i$.
The marginalized likelihood for any data point can be found by multiplying these Gaussians and integrating.
The math is in \cite{HPWL}; the result is
\begin{align}
    \ln p(\vec{x}_i\given \vec{\hat{u}}_D,a) &= -\frac{1}{2}\,[\vec{x}_i - a\,\vec{\hat{u}}_D]^\top\,(\mat{C}_i + \mat{\Lambda})^{-1}\,[\vec{x}_i - a\,\vec{\hat{u}}_D] + \text{constant} ~.
\end{align}
Since the covariance $\mat{\Lambda}$ is zero in one direction, and infinite in all others, this quadratic form depends only on the displacement in the $\vec{\hat{u}}_D$ direction.
That is
\begin{align}
    \ln p(\vec{x}_i\given \vec{\hat{u}}_D,a) &= -\frac{1}{2}\,\frac{(\vec{x}_i^\top\,\vec{\hat{u}}_D - a)^2}{\vec{\hat{u}}_D^\top\,\mat{C}_i^{-1}\,\vec{\hat{u}}_D} + \text{constant} ~,
\end{align}
where the numerator involves the scalar projection of the data point onto the normal direction, and the denominator is the scalar projection of covariance matrix.
Now you can combine these individual data-point likelihoods with any priors you like and do Bayesian inference.
Or you can just maximize the combined likelihood.

HOGG SAYS: Technically, the ``constant'' isn't exactly constant. But it involves an infinity, so I am concerned.

\paragraph{Frequentist approach}
Frequentists don't like to put measures on anything they don't need to; in frequentism we deal with latent parameters by maximizing or profiling.
So our goal is to produce, for each data point $\vec{x}_i$, a profile likelihood, which is the likelihood of the maximum-likelihood point that lies on the hyperplane.
It turns out that this point is the ``closest'' point on the hyperplane, but closest under the definition of distance in which the data point's inverse covariance matrix $\mat{C}_i^{-1}$ is treated as the metric of the space.
The profile likelihood is then just the likelihood for the ``true'' position $\vec{\tilde{x}}$ evaluated at that ``closest'' point.

HOGG SAYS: Do math.

\raggedright
\bibliographystyle{plain}
\bibliography{fitting}

\end{document}
